\section*{Capítulo \ref{ch:g3:continents-and-oceans} --- El globo terráqueo: continentes y océanos}

\begin{solution}{exo:g3:continents-and-oceans:1}
Asia, África, América del Norte, América del Sur, la Antártida, Europa y
Australia. El más grande es Asia; el más pequeño, Australia.
\end{solution}

\begin{solution}{exo:g3:continents-and-oceans:2}
El Pacífico, el Atlántico, el Índico, el Austral y el Ártico. El más grande
es el \omterm{def:g2:water-on-earth:ocean}{océano} Pacífico; el más pequeño, el \omterm{def:g2:water-on-earth:ocean}{océano} Ártico.
\end{solution}

\begin{solution}{exo:g3:continents-and-oceans:3}
El \omterm{def:g3:continents-and-oceans:poles}{Polo Norte} está arriba en un \omterm{def:g3:continents-and-oceans:globe}{globo terráqueo}; el \omterm{def:g3:continents-and-oceans:poles}{Polo Sur} está abajo; el
\omterm{def:g3:continents-and-oceans:equator}{ecuador} da toda la vuelta por el centro.
\end{solution}

\begin{solution}{exo:g3:continents-and-oceans:4}
Europa: norte. Australia: sur. La Antártida: sur. América del Norte: norte.
\end{solution}

\begin{solution}{exo:g3:continents-and-oceans:5}
El \omterm{def:g2:water-on-earth:ocean}{océano} Ártico rodea el \omterm{def:g3:continents-and-oceans:poles}{Polo Norte}. El \omterm{def:g2:water-on-earth:ocean}{océano} Austral rodea la Antártida.
\end{solution}

\begin{solution}{exo:g3:continents-and-oceans:6}
De África a América del Sur: el \omterm{def:g2:water-on-earth:ocean}{océano} Atlántico. De Australia, hacia el
este, a América del Sur: el \omterm{def:g2:water-on-earth:ocean}{océano} Pacífico. De África a Australia: el
\omterm{def:g2:water-on-earth:ocean}{océano} Índico.
\end{solution}

\begin{solution}{exo:g3:continents-and-oceans:7}
Sam vive cerca del \omterm{def:g3:continents-and-oceans:equator}{ecuador}, donde hace calor todo el año. Bo vive cerca de
un \omterm{def:g3:continents-and-oceans:poles}{polo}, donde hace frío todo el año.
\end{solution}

\begin{solution}{exo:g3:continents-and-oceans:8}
$10 - 7 = 3$: su dedo cae en la tierra 3 veces. Es lo que cabe esperar: unos
71 cuadrados de cada 100 de la Tierra son agua, así que unas 7 veces de
cada 10 el dedo se para en el \omterm{def:g2:water-on-earth:ocean}{mar}. (Con solo 10 giros, también podría salir
6 u 8: la cuenta cambia de una partida a otra.)
\end{solution}

\begin{solution}{exo:g3:continents-and-oceans:9}
Para dibujar el \omterm{def:g3:continents-and-oceans:world-map}{mapamundi}, el \omterm{def:g3:continents-and-oceans:globe}{globo} se cortó por una línea que pasa por el
centro del \omterm{def:g2:water-on-earth:ocean}{océano} Pacífico y se aplanó. Una mitad del Pacífico queda en el
borde izquierdo del \omterm{def:g2:drawing-to-map:map}{mapa} y la otra mitad en el borde derecho. En la Tierra
real solo hay un \omterm{def:g2:water-on-earth:ocean}{océano} Pacífico: los dos bordes del \omterm{def:g2:drawing-to-map:map}{mapa} son la misma
línea en el \omterm{def:g3:continents-and-oceans:globe}{globo}.
\end{solution}

\begin{solution}{exo:g3:continents-and-oceans:10}
Tom se equivoca. Del más grande al más pequeño, los \omterm{def:g3:continents-and-oceans:continent}{continentes} son Asia,
África, América del Norte, América del Sur, la Antártida, Europa y
Australia: la Antártida es solo el quinto, mucho más pequeño que Asia. El
\omterm{def:g3:continents-and-oceans:world-map}{mapamundi} le engaña porque un mapa plano estira mucho la tierra cerca de
los \omterm{def:g3:continents-and-oceans:poles}{polos}. La Antártida, que está alrededor del \omterm{def:g3:continents-and-oceans:poles}{Polo Sur}, queda estirada en
una larga franja a lo largo de todo el borde de abajo del \omterm{def:g2:drawing-to-map:map}{mapa}. En el \omterm{def:g3:continents-and-oceans:globe}{globo
terráqueo} es una masa de tierra redonda, y mucho más pequeña que Asia.
\end{solution}
